%! $n$th Roots \& Rational Exponents — Unit 5, Lesson 5.1 — Homework
\documentclass[10pt]{article}
\usepackage{algebra2-article}
\usepackage{algebra2-boxes}
\newcommand{\TallMath}[1]{$\displaystyle #1\rule[-1.4em]{0pt}{3.2em}$}

\begin{document}

\pageheader{Unit 5, Lesson 5.1}{Homework}
\namedateperiod

\vspace{0.10in}

\begin{notesbox}{Practice}
Assume all variables represent \textbf{positive} real numbers except where a problem says otherwise.
No calculator.

\begin{enumerate}[leftmargin=1.9em, itemsep=11pt]
  \item \textbf{Translate both directions.} Fill in every empty cell.
    \begin{center}
    \renewcommand{\arraystretch}{1.9}
    \begin{tabularx}{\linewidth}{@{}X X | X X@{}}
    \hline
    \textbf{Radical form} & \textbf{Rational exponent} & \textbf{Rational exponent} & \textbf{Radical form} \\
    \hline
    $\sqrt[3]{x^{4}}$ & \blank{2.2cm} & $m^{1/5}$   & \blank{2.4cm} \\
    $\sqrt{y^{7}}$    & \blank{2.2cm} & $p^{5/6}$   & \blank{2.4cm} \\
    $\left(\sqrt[4]{w}\right)^{5}$ & \blank{2.2cm} & $(2k)^{2/3}$ & \blank{2.4cm} \\
    \hline
    \end{tabularx}
    \end{center}

  \item \textbf{Evaluate.} Take the root first.
    \begin{center}
    \renewcommand{\arraystretch}{1.9}
    \begin{tabularx}{\linewidth}{@{}X X X X@{}}
    (a) \; $49^{1/2}=$ \blank{1.7cm} & (b) \; $8^{2/3}=$ \blank{1.7cm} &
    (c) \; $81^{3/4}=$ \blank{1.7cm} & (d) \; $32^{2/5}=$ \blank{1.7cm} \\
    (e) \; $100^{-1/2}=$ \blank{1.7cm} & (f) \; $64^{-2/3}=$ \blank{1.7cm} &
    (g) \; $(-125)^{1/3}=$ \blank{1.7cm} & (h) \; $\left(\dfrac{16}{81}\right)^{3/4}=$ \blank{1.7cm} \\
    \end{tabularx}
    \end{center}

  \item \textbf{Simplify with the exponent properties.} Leave every answer with a single
        \emph{positive} rational exponent.
    \begin{center}
    \renewcommand{\arraystretch}{2.0}
    \begin{tabularx}{\linewidth}{@{}X X@{}}
    (a) \; $x^{2/3}\cdot x^{1/6} = $ \blank{2.6cm} &
    (b) \; $\dfrac{c^{7/4}}{c^{1/2}} = $ \blank{2.6cm} \\
    (c) \; $\left(t^{4/5}\right)^{5/2} = $ \blank{2.6cm} &
    (d) \; $\sqrt[3]{x}\cdot\sqrt{x} = $ \blank{2.6cm} \\
    (e) \; $\dfrac{\sqrt[4]{n^{3}}}{\sqrt[8]{n}} = $ \blank{2.6cm} &
    (f) \; $\left(27x^{9}\right)^{2/3} = $ \blank{2.6cm} \\
    \end{tabularx}
    \end{center}

  \item \textbf{The absolute-value rule.} Here $x$ may be \emph{any} real number.
    \begin{center}
    \renewcommand{\arraystretch}{1.8}
    \begin{tabular}{@{}l l l@{}}
    (a) \; $\sqrt{x^{2}}=$ \blank{1.8cm} & (b) \; $\sqrt[3]{x^{3}}=$ \blank{1.8cm} &
    (c) \; $\sqrt[4]{x^{4}}=$ \blank{1.8cm} \\
    \end{tabular}
    \end{center}
    Explain the pattern: which indices need absolute value bars, and why does the other kind not?
    \writelines{3}

  \item \textbf{Back to domains (Lesson 5.0, in exponent language).} For each function, state
        whether the domain is restricted, and give the domain. Look at the \emph{denominator}.
    \begin{center}
    \renewcommand{\arraystretch}{1.8}
    \begin{tabularx}{\linewidth}{@{}l X X X@{}}
    \hline
    \textbf{Function} & \textbf{Denominator: even or odd?} & \textbf{Restricted?} & \textbf{Domain} \\
    \hline
    $f(x)=x^{1/2}$   & \blank{1.8cm} & \blank{1.6cm} & \blank{2.4cm} \\
    $g(x)=x^{1/3}$   & \blank{1.8cm} & \blank{1.6cm} & \blank{2.4cm} \\
    $h(x)=x^{2/5}$   & \blank{1.8cm} & \blank{1.6cm} & \blank{2.4cm} \\
    $j(x)=(x-3)^{1/4}$ & \blank{1.8cm} & \blank{1.6cm} & \blank{2.4cm} \\
    \hline
    \end{tabularx}
    \end{center}

  \item \textbf{Explain.} Yesterday you learned that an \emph{even index} restricts a radical
        function's domain. Today you learned that an \emph{even denominator} does the same thing.
        Explain why these are not two rules but one. \writelines{3}
\end{enumerate}
\end{notesbox}

\begin{extensionbox}
\textbf{Why a bigger animal does not eat proportionally more.} Biologists model an animal's resting
metabolic rate --- roughly, the calories per day it burns just staying alive --- as a rational
power of its body mass. For a mammal of mass $w$ kilograms,
\[
  M(w) = 70\,w^{3/4} \ \text{ kilocalories per day.}
\]
\begin{center}
\renewcommand{\arraystretch}{1.25}
\begin{tabular}{|c|c|c|c|c|}
\hline
$w$ (kg)        & $1$ & $16$ & $81$ & $256$ \\ \hline
$M(w)$ (kcal/day) & $70$ & & & \\ \hline
\end{tabular}
\end{center}
\begin{enumerate}[label=(\alph*), leftmargin=1.8em, itemsep=5pt]
  \item Complete the table without a calculator. (Each mass was chosen to be a perfect fourth power.)
  \item Write $M(w)$ in radical form. Which of the two equivalent radical forms is easier to compute
        with, and why?
  \item From $w=16$ to $w=256$ the mass is multiplied by $16$. By what factor is the metabolic rate
        multiplied? Show the arithmetic, then say in one sentence what the exponent $\frac{3}{4}$
        --- a number \emph{less} than $1$ --- is doing to the growth.
  \item Solve $M=70w^{3/4}$ for $w$ in terms of $M$. Which power did you raise both sides to, and
        why that one?
  \item Your answer to (d) is the \textbf{inverse} of $M$ --- the same relationship run backwards,
        exactly as in Lesson 5.0. What does it let a biologist do that $M$ itself does not?
\end{enumerate}
\writelines{6}
\end{extensionbox}

\begin{spiralbox}
\textbf{Coming next --- Lesson 5.2: Simplifying Radicals; Adding \& Subtracting.} Today you learned
that $\sqrt[n]{a}=a^{1/n}$, which makes an exponent rule out of every radical rule. Next you will use
two of them constantly --- $\sqrt[n]{ab}=\sqrt[n]{a}\cdot\sqrt[n]{b}$ and
$\sqrt[n]{a/b}=\sqrt[n]{a}\big/\sqrt[n]{b}$ --- to pull perfect $n$th powers out from under a
radical ($\sqrt{72}=6\sqrt{2}$), and then discover that radicals add only when they are
\emph{like radicals}, in exactly the way $3x+5x$ works but $3x+5y$ does not.
\end{spiralbox}

\end{document}
